\scanpage{8}
De plus, se rappelant la definition de
\[
 J_n=\widetilde D_n=\text{ens des sous-objet de }D_n
\]
on a, pour \(\forall n\in\mathbb N\), une application canonique, strictement croissante
\[
 J_n\xrightarrow{\tau_n}\mathbb N \qquad\qquad
 \text{\(\tau_n\), application ``type}
\]
dont l'image est formée \deleted{des} de l'ens
\deleted{\(I_{\leq n}\) des} \(\mathbb N_{\leq n}=\{n'\in\mathbb N\mid n'\leq n\}\).

Notons que chaque \(J_n\) a un plus grand élément, qui est ``\(D_n\) (en tant que \inserted{le} sous-objet maximal de \(D_n\)), et que je note \(d_n\). On a donc \inserted{(par (ii))}
\[
 \tau_n^{-1}(n)=\{d_n\}
\]
\deleted{Si \(n'\leq n\), alors pour tout \(x\in J_n\)}

Soient \(n,n'\in\mathbb N\), \deleted{\(\unsure{x\in J|x\in}\)}\deleted{\(x\in J_n,\ n'=\tau_n(x)\), on} \inserted{avec \(n'\leq n\)}
considérons \deleted{\unsure{via \(x\)|via \((x)\)}} \deleted{on a alors \unsure{une|un} application}
\[
 \begin{aligned}
 \mathrm{Hom}(\deleted{\unsure{O|D}}D_{n'},D_n)&\hookrightarrow \widetilde D_n=J_n\\
 u&\longmapsto u(d_{n'})
 \end{aligned}
\]
qui est une bijection des \(\mathrm{Hom}\) sur l'ens \(\tau_n^{-1}(n')\), et considérons
\[
 \mathrm{Hom}(\deleted{\unsure{D|O}}D_{n'}^{\unsure{4|\ast}},D_n)
 \longrightarrow
 \mathrm{Hom}\bigl(\underset{\substack{\tikz{\draw(0,.08)--(0,0)--(.60,0)--(.60,.08);}\\ J_{n'}}}{\widetilde D_{n'}},\underset{\substack{\tikz{\draw(0,.08)--(0,0)--(.54,0)--(.54,.08);}\\J_n}}{\widetilde D_n}\bigr)
\]
on en déduit une \deleted{bijec} application
\[
 \begin{aligned}
 \underset{\substack{\tikz[x=1em,y=1em,baseline=0pt]{\draw[line width=.35pt] (.12,.56) .. controls (.18,.37) and (-.08,.12) .. (.02,.03) .. controls (.07,-.02) and (.18,-.02) .. (.23,.03); \draw[line width=.35pt] (.43,.59) .. controls (.39,.40) and (.39,.20) .. (.42,.08); \fill (0,.61) circle (.025);}\\ {[\mathrm{Hom}(D_{n'},D_n)]}}}{\tau_n^{-1}(n')}
 &\overset{\mathrm{def}}{=}
 \underset{\substack{\tikz{\draw(0,.07)--(0,0)--(.86,0)--(.86,.07);}\\\text{ens des}\\\text{éléments}\\\text{de \(J_n\) de}\\\text{type \(n'\)}}}{J_{n',n}\deleted{\unsure{(n')|(n)}}}
 \xrightarrow{\rho_{n,n'}}\mathrm{Hom}(J_{n'},J_n)\\
 &\hspace{7em}\text{ou}\quad \xi\longmapsto u_{\xi}^{n}
 \end{aligned}
\]
qui est injective, plus précisément on a la \unsure{commutativité|commutativite}
\diagram{\begin{tikzcd}[ampersand replacement=\&,column sep=large,execute at end picture={
 \draw ([xshift=.15em,yshift=-1em]\tikzcdmatrixname-1-1.south east)
 ..controls +(-1.1em,.6em) and +(-2.1em,0)..
 node[pos=.95,above,inner sep=.5pt,fill=white] {\(\scriptstyle\text{inclusion}\)}
 ([xshift=.45em,yshift=-1.45em]\tikzcdmatrixname-1-2.south)
 ..controls +(1.9em,-.04em) and +(-1.0em,-.60em)..([xshift=-1em,yshift=-.75em]\tikzcdmatrixname-1-3.south west);}]
 J_{n,n'} \arrow[r,"\rho_{n,n'}"]
 \& \mathrm{Hom}(J_{n'},J_n) \arrow[r,"u\mapsto u(d_{n'})"] \& J_n
\end{tikzcd}}
% archive sheet number: 7
